
\usepackage[utf8]{inputenc}
\usepackage[english]{babel}
\usepackage[T1]{fontenc}
\usepackage[left=2.5cm,right=2.5cm,top=3cm,bottom=3cm]{geometry}
\usepackage[toc,page]{appendix}
\usepackage{amsfonts,amssymb,amsthm,graphicx,dsfont,xcolor,stmaryrd,subcaption}
\usepackage{amsmath, mathtools}
\usepackage{enumitem}
\usepackage{hyperref}
\usepackage{float} % for the figures
\usepackage{multirow}
% \usepackage{pifont} % usefull for \ding{}
\usepackage[normalem]{ulem} % usefull for underlying
\usepackage{array, booktabs}

\usepackage{mathabx}


% \pdfoutput=1


%%%%%%%%%%%%%%%%%%%%%%%%%%%%%%%%%%%%%%%%%
\DeclareMathOperator{\e}{e}

\DeclareMathOperator{\Proj}{\mathrm{Proj}}
\DeclareMathOperator{\dist}{\mathrm{dist}}
\DeclareMathOperator{\card}{\mathrm{card}}
\DeclareMathOperator{\sign}{\mathrm{sign}}
\DeclareMathOperator{\Conv}{\mathrm{Conv}}


\DeclareMathOperator{\A}{\mathcal{A}}
% \DeclareMathOperator{\L}{\mathcal{L}}
\DeclareMathOperator{\F}{\mathcal{F}}
\DeclareMathOperator{\N}{\mathcal{N}}

% \DeclareMathOperator{\U}{\mathcal{U}}

\DeclareMathOperator{\R}{\mathds{R}}
% \DeclareMathOperator{\R^d}{\R^d}
\DeclareMathOperator{\Integer}{\mathds{N}}
\DeclareMathOperator{\V}{\mathds{V}\mathrm{ar}}
\DeclareMathOperator{\E}{\mathds{E}}
\DeclareMathOperator{\Cov}{\mathds{C}\mathrm{ov}}
\DeclareMathOperator{\Corr}{\mathds{C}\mathrm{orr}}
\DeclareMathOperator{\Prob}{\mathds{P}}
\DeclareMathOperator{\Law}{\mathcal{L}}
\DeclareMathOperator{\Tr}{Tr}

\DeclareMathOperator*{\argmax}{arg\,max}
\DeclareMathOperator*{\argmin}{arg\,min}

\DeclareMathOperator{\1}{\mathds{1}}

\DeclareMathOperator{\C1}{\mathcal{C}^1}
\DeclareMathOperator{\C1Lip}{\mathcal{C}^1_{_{Lip}}}
\DeclareMathOperator\supp{supp}


\DeclareMathOperator{\Distortion}{\mathcal{Q}_{2,N}}

\newcommand\indep{\protect\mathpalette{\protect\indeT}{\perp}}
\def\indeT#1#2{\mathrel{\rlap{$#1#2$}\mkern2mu{#1#2}}}

\newcommand{\vertiii}[1]{{\left\vert\kern-0.25ex\left\vert\kern-0.25ex\left\vert #1
\right\vert\kern-0.25ex\right\vert\kern-0.25ex\right\vert}}

\theoremstyle{plain}
\newtheorem{theorem}{Theorem}[section]
\newtheorem{lemme}[theorem]{Lemma}
\newtheorem{proposition}[theorem]{Proposition}
\newtheorem{cor}[theorem]{Corollary}

\theoremstyle{definition}
\newtheorem{definition}[theorem]{Definition}
\newtheorem{conj}[theorem]{Conjecture}
\newtheorem{example}[theorem]{Example}

% \theoremstyle{remark}
\newtheorem{remark}[theorem]{\noindent \textbf{Remark}}
\newtheorem{remarks}[theorem]{\noindent \textbf{Remarks}}
\newtheorem*{note}{Note}



%%%%%%%%%%%%%%%%%%%%%%%%%%%%%%%%%%%%%%%%%

\numberwithin{equation}{section}

\renewenvironment{abstract}
 {\small
	\begin{center}
		\bfseries \abstractname\vspace{-.5em}\vspace{0pt}
	\end{center}
	\list{}{%
	\setlength{\leftmargin}{20mm}% <---------- CHANGE HERE
	\setlength{\rightmargin}{\leftmargin}%
 }%
 \item\relax}
 {\endlist}




% Keywords command
\providecommand{\keywords}[1]
{
  {
  \small
  \textbf{\textit{Keywords---}} #1
  }
}



\newcommand*\samethanks[1][\value{footnote}]{\footnotemark[#1]}
